%% kapa: exam=1 round=28 year=2017 period=1 form=A subject=civil_law subject_ko=민법 seq=31 number=31
%% source: Q-Net 감정평가사 기출문제 (cst00309) artlSeq=5199603 file=extracted/1008/5199603/2017년 제28회 감정평가사 1차시험 문제지/1교시 A형.pdf
%% conversion: kordoc 4.24.51 -> tools/gen1.py (auto); review=proofread-sonnet; flags=none
\begin{question}{31}
부동산 실권리자명의 등기에 관한 법률상 명의신탁에 대한 설명으로 옳지 않은 것은? (다툼이 있으면 판례에 따름)
\begin{choices}
  \item 무효인 명의신탁등기가 행하여진 후 신탁자와 수탁자가 혼인한 경우, 조세포탈 등의 목적이 없더라도 그 명의신탁등기는 유효로 인정될 수 없다.
  \item 채무변제를 담보하기 위해 채권자 명의로 부동산에 관한 소유권이전등기를 하기로 하는 약정은 명의신탁약정에 해당하지 않는다.
  \item 무효인 명의신탁약정에 기하여 타인 명의의 등기가 마쳐졌다는 이유만으로 그것이 당연히 불법원인급여에 해당한다고 볼 수 없다.
  \item 조세포탈 등의 목적 없이 종교단체 명의로 그 산하조직이 보유한 부동산의 소유권을 등기한 경우, 그 단체와 조직 간의 명의신탁약정은 유효하다.
  \item 신탁자는 명의신탁약정의 무효로서 수탁자로부터 소유권이전등기를 받은 제3자에게 그의 선의ㆍ악의 여부를 불문하고 대항하지 못한다.
\end{choices}
\end{question}
\answer{1}
